%=======================================================================
\chapter{Benchmark des modèles de Machine Learning}

\section{Modèles évalués}

L'objectif principal de cette phase est d'identifier la combinaison la plus performante entre
les ensembles de caractéristiques construits lors de la Phase~2 et différents algorithmes de
Machine Learning. Afin d'obtenir une évaluation exhaustive, huit modèles de classification
représentant différentes familles d'apprentissage supervisé ont été sélectionnés.

\subsection{Support Vector Machine (SVM)}

Le Support Vector Machine est un classificateur fondé sur la recherche d'un hyperplan optimal
maximisant la marge entre les différentes classes. Grâce à l'utilisation de fonctions noyaux,
il est capable de traiter efficacement des données non linéaires et de grande dimension.

Le SVM est particulièrement adapté aux problèmes de classification médicale impliquant un
nombre élevé de caractéristiques, ce qui en fait un candidat naturel pour les ensembles
fusionnés étudiés dans ce projet.

\subsection{K-Nearest Neighbors (KNN)}

Le classificateur KNN repose sur un mécanisme de vote majoritaire parmi les $k$ voisins les plus
proches d'un échantillon.

Sa simplicité et son absence d'hypothèses fortes sur la distribution des données en font une
référence classique pour les tâches de classification multiclasses.

\subsection{Random Forest (RF)}

Random Forest est une méthode d'apprentissage par ensemble basée sur la construction de
plusieurs arbres de décision indépendants.

La prédiction finale est obtenue par vote majoritaire des arbres individuels. Cette approche
offre généralement une bonne robustesse face au bruit et aux variables redondantes.

\subsection{Extra Trees (ET)}

L'algorithme Extra Trees constitue une variante plus aléatoire de Random Forest.

L'introduction d'une randomisation supplémentaire lors de la construction des arbres permet
souvent de réduire la variance du modèle tout en conservant une bonne capacité de
généralisation.

\subsection{Logistic Regression (LR)}

La Régression Logistique est un modèle linéaire probabiliste largement utilisé comme
référence dans les études comparatives.

Bien qu'elle soit plus simple que les autres modèles évalués, elle permet d'analyser la
séparabilité intrinsèque des ensembles de caractéristiques.

\subsection{Multi-Layer Perceptron (MLP)}

Le Multi-Layer Perceptron est un réseau de neurones artificiel composé de plusieurs couches
entièrement connectées.

Grâce à ses fonctions d'activation non linéaires, il est capable de modéliser des frontières
de décision complexes et d'apprendre des relations non linéaires entre les caractéristiques.

\subsection{XGBoost}

XGBoost (\emph{Extreme Gradient Boosting}) est une méthode de boosting basée sur l'agrégation
séquentielle d'arbres de décision.

Son efficacité, sa robustesse et sa capacité à gérer des données de grande dimension en font
l'un des algorithmes les plus performants dans de nombreux problèmes de classification.

\subsection{LightGBM}

LightGBM est une implémentation optimisée du Gradient Boosting développée pour offrir des
temps d'entraînement réduits et une meilleure gestion des grands ensembles de données.

Sa structure efficace lui permet de construire rapidement des modèles performants tout en
limitant la consommation mémoire.

% -----------------------------------------------------------------------
\section{Stratégie d'évaluation}

Afin d'assurer une comparaison équitable entre les différents modèles, une procédure
d'évaluation rigoureuse a été mise en place.

\subsection{Optimisation des hyperparamètres par GridSearchCV}

Chaque algorithme possède plusieurs hyperparamètres susceptibles d'influencer fortement ses
performances.

Afin d'identifier automatiquement les meilleures configurations, une recherche par grille
(\emph{Grid Search}) a été réalisée à l'aide de la méthode \textbf{GridSearchCV} de
Scikit-Learn.

Cette procédure consiste à~:
\begin{enumerate}[leftmargin=1.4em,itemsep=2pt]
  \item générer toutes les combinaisons possibles d'hyperparamètres~;
  \item entraîner le modèle pour chaque combinaison~;
  \item évaluer les performances via validation croisée~;
  \item retenir la configuration obtenant le meilleur score moyen.
\end{enumerate}

Cette approche garantit que chaque modèle est évalué dans des conditions optimales.

\subsection{Validation croisée stratifiée}

Pour réduire l'impact des variations liées à la partition des données, une validation croisée
stratifiée à \cvFolds{} plis (\emph{Stratified \cvFolds{}-Fold Cross Validation}) a été
utilisée sur l'ensemble d'apprentissage (\NTrain{} images).

Le principe consiste à~:
\begin{itemize}[leftmargin=1.4em,itemsep=2pt]
  \item diviser le jeu d'apprentissage en cinq sous-ensembles~;
  \item utiliser successivement quatre plis pour l'entraînement et un pli pour la validation~;
  \item répéter l'opération cinq fois~;
  \item calculer la moyenne des performances obtenues.
\end{itemize}

L'utilisation d'une stratégie stratifiée permet de conserver la proportion des classes dans
chaque pli.

\subsection{Métrique d'évaluation : F1-Macro}

La métrique principale retenue dans ce projet est le \textbf{F1-Macro Score}, utilisée à la
fois pour la sélection des hyperparamètres (score de validation croisée) et pour le
classement final sur l'ensemble de test (\NTest{} images).

Cette mesure combine la précision (\emph{Precision}) et le rappel (\emph{Recall})~:
\begin{equation}
  F1 = 2 \times \frac{\mathrm{Precision} \times \mathrm{Recall}}
                    {\mathrm{Precision} + \mathrm{Recall}}
\end{equation}

Dans sa version Macro, le score est calculé indépendamment pour chaque classe puis moyenné.

Cette approche est particulièrement adaptée aux problèmes multiclasses car elle accorde le même
poids à chaque catégorie de tumeur.

Le F1-Macro constitue ainsi un indicateur robuste de la performance globale des modèles.

% -----------------------------------------------------------------------
\section{Résultats détaillés}

Chaque combinaison «~Ensemble de caractéristiques $\times$ Modèle~» a été évaluée selon la
procédure décrite précédemment.

Au total, huit ensembles de caractéristiques ont été testés sur huit algorithmes de Machine
Learning, soit \textbf{64 expériences} distinctes. Les résultats complets sont présentés dans
l'Annexe~\ref{annex:benchmark-results}. Ils permettent d'analyser~:
\begin{itemize}[leftmargin=1.4em,itemsep=2pt]
  \item l'influence du choix des caractéristiques~;
  \item l'impact de la fusion des descripteurs~;
  \item les différences de comportement entre les algorithmes~;
  \item la stabilité des performances observées.
\end{itemize}

Les premières observations montrent que les ensembles fusionnés obtiennent généralement de
meilleurs résultats que les ensembles individuels, confirmant l'intérêt de combiner
plusieurs sources d'information visuelle.

Par ailleurs, les modèles non linéaires tels que SVM, KNN, LightGBM et XGBoost bénéficient
davantage de la richesse descriptive apportée par les ensembles de caractéristiques fusionnés.

Les Figures~\ref{fig:p2-A} à \ref{fig:p2-Full} présentent les performances détaillées
(F1-Macro, précision, rappel et matrices de confusion) pour chaque ensemble de
caractéristiques.

\begin{figure}[htbp]
  \centering
  \includegraphics[width=\textwidth]{p2_A_stats}
  \caption{Comparaison des modèles sur l'ensemble A (Statistiques).}
  \label{fig:p2-A}
\end{figure}

\begin{figure}[htbp]
  \centering
  \includegraphics[width=\textwidth]{p2_B_tex}
  \caption{Comparaison des modèles sur l'ensemble B (Texture optimisée).}
  \label{fig:p2-B}
\end{figure}

\begin{figure}[htbp]
  \centering
  \includegraphics[width=\textwidth]{p2_C_dwt}
  \caption{Comparaison des modèles sur l'ensemble C (DWT optimisé).}
  \label{fig:p2-C}
\end{figure}

\begin{figure}[htbp]
  \centering
  \includegraphics[width=\textwidth]{p2_D_hog}
  \caption{Comparaison des modèles sur l'ensemble D (HOG).}
  \label{fig:p2-D}
\end{figure}

\begin{figure}[htbp]
  \centering
  \includegraphics[width=\textwidth]{p2_ApB}
  \caption{Comparaison des modèles sur l'ensemble fusionné A+B.}
  \label{fig:p2-ApB}
\end{figure}

\begin{figure}[htbp]
  \centering
  \includegraphics[width=\textwidth]{p2_BpC}
  \caption{Comparaison des modèles sur l'ensemble fusionné B+C.}
  \label{fig:p2-BpC}
\end{figure}

\begin{figure}[htbp]
  \centering
  \includegraphics[width=\textwidth]{p2_ApBpC}
  \caption{Comparaison des modèles sur l'ensemble fusionné A+B+C.}
  \label{fig:p2-ApBpC}
\end{figure}

\begin{figure}[htbp]
  \centering
  \includegraphics[width=\textwidth]{p2_Full}
  \caption{Comparaison des modèles sur l'ensemble Full (fusion complète).}
  \label{fig:p2-Full}
\end{figure}

% -----------------------------------------------------------------------
\section{Leaderboard global}

Afin d'identifier les meilleures configurations du projet, l'ensemble des résultats a été
regroupé dans un classement global basé sur le F1-Macro obtenu sur l'ensemble de test.

Le Tableau~\ref{tab:leaderboard-top10} présente les dix combinaisons les plus performantes
(source~: \texttt{phase2\_benchmark\_results.csv}).

\begin{table}[htbp]
  \centering
  \caption{Leaderboard global — Top 10 (F1-Macro sur l'ensemble de test).}
  \label{tab:leaderboard-top10}
  \begin{tabular}{@{}cllc@{}}
    \toprule
    \thd{Rang} & \thd{Modèle} & \thd{Ensemble} & \thd{F1-Macro} \\
    \midrule
    1  & SVM         & Full(opt) & 0,920 \\
    2  & SVM         & D (HOG)   & 0,903 \\
    3  & KNN         & Full(opt) & 0,900 \\
    4  & KNN         & D (HOG)   & 0,897 \\
    5  & LightGBM    & Full(opt) & 0,894 \\
    6  & MLP         & Full(opt) & 0,884 \\
    7  & XGBoost     & Full(opt) & 0,880 \\
    8  & Extra Trees & Full(opt) & 0,877 \\
    9  & MLP         & D (HOG)   & 0,877 \\
    10 & LR          & Full(opt) & 0,870 \\
    \bottomrule
  \end{tabular}
\end{table}

La Figure~\ref{fig:p2-leaderboard} visualise ce classement ; la
Figure~\ref{fig:p2-dashboard} offre une vue d'ensemble des performances sur l'ensemble des
64 expériences.

\begin{figure}[htbp]
  \centering
  \includegraphics[width=0.92\textwidth]{p2_leaderboard}
  \caption{Leaderboard global des combinaisons modèle $\times$ ensemble de caractéristiques.}
  \label{fig:p2-leaderboard}
\end{figure}

\begin{figure}[htbp]
  \centering
  \includegraphics[width=\textwidth]{p2_overall_dashboard}
  \caption{Tableau de bord global du benchmark Phase~2.}
  \label{fig:p2-dashboard}
\end{figure}

L'analyse du leaderboard met en évidence plusieurs tendances importantes.

Premièrement, les ensembles de caractéristiques fusionnés dominent largement les ensembles
individuels, ce qui confirme l'existence d'une forte complémentarité entre les informations
statistiques, texturales, fréquentielles et structurelles. L'ensemble \textbf{Full(opt)}
apparaît dans sept des dix meilleures configurations.

Deuxièmement, le modèle \textbf{SVM} obtient les meilleures performances globales
(F1-Macro = \bestFone{} sur Full(opt)). Ce résultat confirme sa capacité à exploiter
efficacement des espaces de caractéristiques de grande dimension tout en maintenant une
excellente capacité de généralisation.

Troisièmement, les modèles \textbf{KNN} et \textbf{LightGBM} occupent également les premières
positions du classement, démontrant que les informations extraites possèdent une forte
capacité discriminante indépendamment du paradigme d'apprentissage utilisé.

Enfin, l'ensemble \textbf{D (HOG)} seul atteint des performances proches de Full(opt) avec
SVM et KNN (rangs 2 et 4), soulignant l'apport majeur des descripteurs de contours dans ce
problème de classification.

Ces observations seront approfondies dans le chapitre suivant à travers une validation
statistique complète visant à évaluer la robustesse et la significativité des différences
observées entre les modèles.
